\documentclass[10pt,a4paper]{article}
\usepackage[utf8]{inputenc}
\usepackage[a4paper,left=3cm,right=2cm]{geometry}
\usepackage{amsmath}
\usepackage[thmmarks, amsmath, thref]{ntheorem}
\usepackage{amsfonts}
\usepackage{amssymb}
\usepackage{fancyhdr}
\pagestyle{fancy}
\usepackage{anysize}
\usepackage{graphicx}
\usepackage{float}
\usepackage{hyperref}
\usepackage{multirow}
\usepackage{enumitem}
\newtheorem{theorem}{Theorem}
\theoremsymbol{\ensuremath{\blacksquare}}
\newtheorem*{solution}{Solución:}
\usepackage{listings}
\usepackage{xcolor}

% Margenes y formalidades

\textwidth 6.25in % Margen Derecho y ancho del texto
\textheight 8.5in
\oddsidemargin 0in % Margen Izquierdo
\headheight 0in % Largo del texto
\leftmargin 1in
\parindent 0em % Salto entre la indentación (Sangría)https://www.overleaf.com/4519127669sdpptkhvsbfn
\topmargin -0.5in % Margen Superior
\parskip 0.5ex % Salto entre parrafos
\baselineskip 1pt

%\lhead{\includegraphics[scale=0.2]{logo.png}} % texto izquierda de la cabecera
%\chead{}                                      % texto centro de la cabecera
%\rhead{\includegraphics[scale=0.2]{FIG/logo-mba-2016.png}} % número de página a la derecha

\renewcommand{\headrulewidth}{0.4pt} % grosor de la línea de la cabecera
\renewcommand{\footrulewidth}{0.4pt} % grosor de la línea del pie

\begin{document}



\pagebreak
\vspace*{4mm} 

\begin{center}
\begin{huge}
\textbf{\\Ejercicios Econometría 2}
\end{Large}\\
%{\large 2023-2}\\
\end{center}

\begin{enumerate}

  \item  Se cuenta con la siguiente información para 5 individuos, donde $Y$ es la variable dependiente y $X_1$, $X_2$ son variables explicativas:

\[
\begin{array}{c|ccc}
\hline
\text{Individuo} & X_1 & X_2 & Y \\
\hline
1 & 4 & 1 & 5 \\
2 & 3 & 2 & 7 \\
3 & 5 & 3 & 10 \\
4 & 7 & 4 & 13 \\
5 & 6 & 4 & 12 \\
\hline
\end{array}
\]

\begin{enumerate}
  \item Estime el modelo $Y = \beta_0 + \beta_1\,X_1 + \varepsilon$, ignorando la variable $X_2$. Calcule los coeficientes, el $R^2$, el $R^2$ ajustado y las sumas de cuadrados (SST, SSR y SSE).
  \item Discuta la posibilidad de sesgo al omitir $X_2$.
  \item Estime el modelo $Y = \beta_0 + \beta_1 X_1 + \beta_2 X_2 + \varepsilon$, calcule el nuevo $R^2$ (y las sumas de cuadrados) y compare con el modelo simple. Hint: 
  \[
(X^TX)^{-1}
= \frac{1}{95}
\begin{bmatrix}
281 & -72 & 35\\
-72 & 34 & -35\\
35 & -35 & 50
\end{bmatrix}.
\]
\end{enumerate}

\textbf{Soluci\'on}


Deseamos estimar 
\[
Y = \beta_0 + \beta_1 X_1 + \varepsilon,
\]
ignorando la variable $X_2$.

\paragraph{Media de $X_1$ y de $Y$.}

\[
\bar{X}_1
= \frac{4+3+5+7+6}{5}
= \frac{25}{5}
= 5,
\quad
\bar{Y}
= \frac{5+7+10+13+12}{5}
= \frac{47}{5}
= 9.4.
\]

\paragraph{Varianza, covarianza y c\'alculo de $\beta_1$ y $\beta_0$.}

Para la regresión simple:
\[
\beta_1
= \frac{\mathrm{Cov}(X_1,Y)}{\mathrm{Var}(X_1)},
\quad
\beta_0
= \bar{Y} - \beta_1\,\bar{X}_1.
\]

Construyamos la tabla con $(X_1-\bar{X}_1)$, $(Y-\bar{Y})$, y sus productos.

\[
\begin{array}{c|cc|cc|ccc}
\hline
\text{Ind} & X_1 & Y 
& X_1-\bar{X}_1 & Y-\bar{Y}
& (X_1 - \bar{X}_1)^2 & (Y - \bar{Y})^2
& (X_1-\bar{X}_1)(Y-\bar{Y})\\
\hline
1 & 4 & 5 
& 4-5=-1
& 5-9.4=-4.4
& 1
& 19.36
& (+4.4) \\[6pt]
2 & 3 & 7
& 3-5=-2
& 7-9.4=-2.4
& 4
& 5.76
& (+4.8) \\[6pt]
3 & 5 & 10
& 5-5=0
& 10-9.4=0.6
& 0
& 0.36
& 0 \\[6pt]
4 & 7 & 13
& 7-5=2
& 13-9.4=3.6
& 4
& 12.96
& 7.2 \\[6pt]
5 & 6 & 12
& 6-5=1
& 12-9.4=2.6
& 1
& 6.76
& 2.6 \\[6pt]
\hline
\text{Suma} & & & & 
& 10
& 45.20
& 19.0 \\[6pt]
\hline
\end{array}
\]

De la última fila:
\[
\sum (X_1-\bar{X}_1)^2 = 10,
\quad
\sum (Y - \bar{Y})^2 = 45.20,
\quad
\sum (X_1-\bar{X}_1)(Y-\bar{Y}) = 19.0.
\]
Entonces
\[
\mathrm{Var}(X_1)
= \frac{1}{n-1}\sum (X_1-\bar{X}_1)^2
= \frac{10}{4}
= 2.5,
\quad
\mathrm{Cov}(X_1,Y)
= \frac{19.0}{4}
= 4.75.
\]
\[
\beta_1
= \frac{4.75}{2.5}
= 1.90,
\quad
\beta_0
= 9.4 - (1.90)(5)
= 9.4 - 9.5
= -0.1.
\]

\[
\boxed{
\widehat{\beta}_0 \approx -0.10,
\quad
\widehat{\beta}_1 \approx 1.90.
}
\]

\paragraph{ C\'alculo de SST, SSE, SSR y $R^2$.}

\[
\text{SST}
= \sum (Y_i - \bar{Y})^2
= 45.20.
\]
El valor ajustado: 
\[
\hat{Y}_i
= -0.1 +1.90\,X_{1i}.
\]
Calculamos los residuos $e_i = Y_i - \hat{Y}_i$:

\[
\begin{array}{c|cc|c|c|c}
\hline
\text{Ind} & X_1 & Y 
& \hat{Y}_i = -0.1 +1.90X_1
& e_i= Y-\hat{Y}_i
& e_i^2
\\\hline
1 & 4 & 5 
& -0.1 +1.90\cdot4 =-0.1 +7.6=7.5
& 5-7.5=-2.5
& 6.25
\\
2 & 3 &7
& -0.1 +1.90\cdot3 =-0.1 +5.7=5.6
& 7-5.6=1.4
& 1.96
\\
3 & 5 &10
& -0.1 +1.90\cdot5 =-0.1 +9.5=9.4
& 10-9.4=0.6
& 0.36
\\
4 & 7 &13
& -0.1 +1.90\cdot7 =-0.1 +13.3=13.2
& 13-13.2=-0.2
& 0.04
\\
5 & 6 &12
& -0.1 +1.90\cdot6 =-0.1 +11.4=11.3
& 12-11.3=0.7
& 0.49
\\
\hline
& & & & \sum e_i^2= & 9.10
\\\hline
\end{array}
\]
\[
\text{SSR} = \sum e_i^2 \approx 9.10,
\quad
\text{SSE} = \text{SST} - \text{SSR} = 45.20 -9.10=36.10.
\]
\[
R^2
= 1 - \frac{\text{SSR}}{\text{SST}}
= 1 - \frac{9.10}{45.20}
\approx 0.7991.
\]
Es decir, el \textbf{modelo simple} explica alrededor de $80\%$ de la variabilidad de $Y$.

\paragraph{$R^2$ ajustado (modelo simple).}

Con $n=5$, $k=1$:
\[
R_{\text{ajustado}}^2
= 1 - \frac{\text{SSR}/(n - k -1)}{\text{SST}/(n-1)}
= 1 - \frac{9.10/3}{45.20/4}
= 1 - \frac{3.0333}{11.30}
\approx 1 -0.2684
= 0.7316.
\]


\[
Y \approx -0.10 +1.90\,X_1,
\quad
R^2 \approx 0.799,
\quad
R_{\text{ajustado}}^2 \approx 0.732.
\]

b) \\
La variable $X_2$ podr\'ia ser relevante. Si $X_2$ est\'a correlacionada con $X_1$ o explica parte de la variabilidad de $Y$, omitirla puede causar que la pendiente estimada $\beta_1$ en la regresión simple sea \emph{sesgada}. Para observar si este es el caso, estimamos ahora el modelo con ambas variables y vemos si $R^2$ crece y si $\beta_1$ se modifica.

c) \\

El modelo completo:
\[
Y = \beta_0 + \beta_1\,X_1 + \beta_2\,X_2 + \varepsilon.
\]
En notaci\'on matricial:
\[
\hat{\boldsymbol{\beta}}
= (\mathbf{X}^\top \mathbf{X})^{-1} \, (\mathbf{X}^\top \mathbf{Y}),
\]
donde
\[
\mathbf{X}
=
\begin{bmatrix}
1 & X_1 & X_2
\end{bmatrix},
\quad
\mathbf{Y}
=
\begin{bmatrix}
Y
\end{bmatrix}.
\]


\[
\mathbf{X}
=
\begin{bmatrix}
1 & 4 & 1\\
1 & 3 & 2\\
1 & 5 & 3\\
1 & 7 & 4\\
1 & 6 & 4
\end{bmatrix},
\quad
\mathbf{Y}
=
\begin{bmatrix}
5\\
7\\
10\\
13\\
12
\end{bmatrix}.
\]



{$\hat{\boldsymbol{\beta}} = (X^TX)^{-1}(\mathbf{X}^\top \mathbf{Y})$}

Sea
\[
\mathbf{X}^\top \mathbf{Y}
=
\begin{bmatrix}
47\\
254\\
149
\end{bmatrix}.
\]

Primero calculamos 
\(\mathrm{adj}(A)\cdot b\),
sin dividir todavía por 95:
\[
\begin{bmatrix}
281 & -72 & 35\\
-72 & 34 & -35\\
35 & -35 & 50
\end{bmatrix}
\begin{bmatrix}
47\\
254\\
149
\end{bmatrix}
=
\begin{bmatrix}
\text{fila1}\\
\text{fila2}\\
\text{fila3}
\end{bmatrix}.
\]
\[
\text{(fila 1)} 
= 281\cdot47 + (-72)\cdot254 + 35\cdot149.
\]
\[
281\cdot47=13207,\quad
(-72)\cdot254=-18288,\quad
35\cdot149=5215.
\]
Suma:
\[
13207 -18288 +5215= 13207 +5215=18422,\quad 18422-18288=134.
\]
\[
\text{(fila 2)}
= (-72)\cdot47 +34\cdot254 +(-35)\cdot149.
\]
\[
(-72)\cdot47=-3384,\quad
34\cdot254=8636,\quad
(-35)\cdot149=-5215.
\]
Suma:
\[
-3384+8636=5252,\quad 5252-5215=37.
\]
\[
\text{(fila 3)}
= 35\cdot47 +(-35)\cdot254 +50\cdot149.
\]
\[
35\cdot47=1645,\quad
(-35)\cdot254=-8890,\quad
50\cdot149=7450.
\]
Suma:
\[
1645-8890 +7450=9095-8890=205.
\]
El vector intermedio es
\[
\begin{bmatrix}
134\\
37\\
205
\end{bmatrix}.
\]
Dividimos por $95$:
\[
\hat{\boldsymbol{\beta}}
=
\frac{1}{95}
\begin{bmatrix}
134\\
37\\
205
\end{bmatrix}
=
\begin{bmatrix}
134/95\\
37/95\\
205/95
\end{bmatrix}
\approx
\begin{bmatrix}
1.41\\
0.39\\
2.16
\end{bmatrix}.
\]
\[
\boxed{
\hat{\beta}_0 \approx 1.41,
\quad
\hat{\beta}_1 \approx 0.39,
\quad
\hat{\beta}_2 \approx 2.16.
}
\]


\paragraph{SSE, SSR, SST.}

La \emph{Suma Total de Cuadrados} (SST) es
\[
\text{SST}
= \sum_{i=1}^5 (\,Y_i - \bar{Y}\,)^2.
\]
En este caso, la SST resulta $45.20$ (véase el modelo simple).  

El valor ajustado es
\[
\hat{Y}_i
= 1.41 + 0.39\,X_{1i} + 2.16\,X_{2i}.
\]
Los residuos $e_i=Y_i-\hat{Y}_i$ dan la \emph{Suma de Cuadrados de Residuos} (SSR). Tras los cálculos numéricos (ver tabla de $e_i^2$), se obtiene
\[
\text{SSR} \approx 0.253.
\]
Entonces
\[
\text{SSE}
= \text{SST} - \text{SSR}
= 45.20 - 0.253
= 44.947.
\]
\[
R^2
= 1-\frac{\text{SSR}}{\text{SST}}
= 1 - \frac{0.253}{45.20}
\approx 0.9944.
\]

\paragraph{$R_{\text{ajustado}}^2$.}

Con $n=5$ y $k=2$ variables explicativas (sin contar la constante), la fórmula es:
\[
R_{\text{ajustado}}^2
= 1
- 
\frac{\text{SSR}/(n-k-1)}{\text{SST}/(n-1)},
\]
es decir
\[
= 1
- 
\frac{0.253/(5-2-1)}{45.20/(5-1)}
= 1
- 
\frac{0.253/2}{45.20/4}.
\]
\[
0.253/2 = 0.1265,
\quad
45.20/4=11.30,
\quad
\frac{0.1265}{11.30}\approx 0.0112,
\]
\[
R_{\text{ajustado}}^2
= 1-0.0112
= 0.9888.
\]
\[
\boxed{
R^2 \approx 0.9944,
\quad
R_{\text{ajustado}}^2 \approx 0.9888.
}
\]


\paragraph{C\'alculos en R.} Usamos \texttt{lm()}:

\begin{lstlisting}
# Datos
X1 <- c(4,3,5,7,6)
X2 <- c(1,2,3,4,4)
Y  <- c(5,7,10,13,12)
datos <- data.frame(X1, X2, Y)

# Modelo simple:
mod_simple <- lm(Y ~ X1, data=datos)
summary(mod_simple)
# (Muestra coeficientes, R^2, etc.)

# Modelo multiple:
mod_multiple <- lm(Y ~ X1 + X2, data=datos)
summary(mod_multiple)
# (Nuevos coeficientes, R^2 mayor, etc.)
\end{lstlisting}





    \item Utilizando una base de datos de 4.137 alumnos universitarios, se estimó la ecuación siguiente de MCO:

    
$$\hat{colgpa} =1,392 -0,0135 hsperc + 0,0014 sat$$

Con $R^2=0,273$ y donde $colgpa$ es el promedio de las calificaciones de los estudiantes universitarios,
$hsperc$ es el percentil en la clase de educación media en la que se graduó (definida de manera que, por
ejemplo, $hsperc= 5$ significa el 5\% superior de su clase), y $sat$ son los puntajes combinados en
matemáticas y habilidades verbales en la prueba de admisión universitaria de los alumnos.
\begin{enumerate} 
\item (8 puntos) ¿Por qué es lógico que el coeficiente en $hsperc$ sea negativo?

{\color{blue} Un valor  bajo en $hsperc$, implica que el estudiante obtuvo alguno de los mejores resultados de su clase. Se espera que los mejores estudiantes del colegio, también tengan buenas calificaciones en la universidad.  }
\item  (8 puntos) ¿Cual es el promedio de calificaciones universitarias predicho cuando $hsperc= 20$ y $sat=900$?

{\color{blue} $$\hat{colgpa} =1,392 -0,0135 \cdot 20 + 0,0014 \cdot 900=2,382$$}


\item (10 puntos) Suponga que dos estudiantes, A y B, se graduaron en el mismo percentil de la enseñanza media, pero el puntaje $sat$ del estudiante A es 140 puntos más alto. ¿Cuál es la diferencia predicha en el promedio
de calificaciones universitarias para estos dos alumnos? 

{\color{blue} Para esto usamos el parámetro estimado para sat 0,0014. Si todo lo demás está constante, entonces la diferencia en las calificaciones de A y B predichas será $0,0014\cdot 140=0,196$.}

\item (10 puntos) Determine el $R^2$ ajustado ($\Bar{R^2}$) de la regresión. Compárelo con el $R^2$ y explique.

{\color{blue}
$SSR/SST=	0,727$\\
$\bar{R^2}=\frac{n-1}{n-k-1}0.727=\frac{4137-1}{4137-2-1}0.727$\\
$\bar{R^2}=0.2726$

Es menor al $R^2$ dado que el $\bar{R^2}$ penaliza por la disminución de grados de de libertad al incluir varios regresores al modelo, aunque su variación es pequeña dado que se tiene una gran cantidad de datos versus la cantidad de variables del modelo.}

\item (20 puntos) Considerando que:

\[(X^\top X)^{-1} =
\begin{bmatrix}
8,5 & 0,19 & -0,11 \\
0,19 & 0,5 & -0,02 \\
-0,11 & -0,02 & 0,04 \\
\end{bmatrix}
\]

$$ \sum_{i=1}^{n} (y_i - \bar{y})^2=6.800$$
Determine la varianza de los estimadores de la regresión MCO.

{\color{blue} 
Sabemos que $SSR/SST = 0,727$ y $SST=6800$, por lo tanto, $SSR = 0,727\cdot 6800=4943,6 $, entonces:
$$\hat{\sigma}^2 = \frac{\sum_{i=1}^{n} \hat{u}_i^2 }{n-k-1}= \frac{SSR}{n - k - 1}=\frac{4943,6}{4137-2-1}=1,1958$$
    
    \[
\text{Var}(\hat{\beta}|X) = \hat{\sigma}^2(X^{\top}X)^{-1}
\]

Solo necesitamos la varianza de los estimadores así que  os quedamos con los elementos de la diagonal de la matriz de var-covar.

$$\text{Var}(\hat{\beta_0}|X)=10,1643\quad 
\text{Var}(\hat{\beta_1}|X)=0,5979
\quad \text{Var}(\hat{\beta_2}|X)=0,0478$$}
\end{enumerate}


\item (Continuación ejercicio 1) Se tienen los siguientes datos de salario, años de educación y experiencia laboral de 5 individuos:
\begin{center}
\begin{tabular}{|c|c|c|c|}
\hline
\textbf{Individuo} & \textbf{Educ} ($X_1$) & \textbf{Exp} ($X_2$)  & \textbf{Salario ($Y$)}\\
\hline
1 & 12 & 14 & 1000 \\
2 & 14 & 15 & 1200 \\
3 & 15 & 17 & 1300 \\
4 & 16 & 15 & 1200 \\
5 & 18 & 16 & 1500 \\
\hline
\end{tabular}
\end{center}

Se hizo el modelo en torno solo a eduación de la forma
$$\hat{Salario}= 115 + 75\cdot educ$$
con $R^2 = 0.85$ y $R^2_{adj} = 0.8$
\begin{enumerate}

    \item (5 puntos) ¿Qué sesgos puede tener el modelo anterior?
    \\
    \textbf{solución}
    \\
    El modelo anterior puede tener sesgo por omisión de variables, ya que la variable experiencia puede ser relevante en el modelo.
    \\

    \item (10 puntos) Luego se realiza un modelo incluyendo la variable Experiencia, por lo que el salario estaría explicado por la educación y la experiencia de la siguiente manera:
    $$\widehat{Salario} = -430 + 60 \cdot Educ + 50 \cdot Exp$$
    Calcule SST, SSE, SSR, $R^2$ y $R^2$ ajustado.
    \\
    \textbf{solución}
    \\
    Acá SST es el mismo valor de 132000. Calcularemos SSE.
    \begin{center}
\begin{tabular}{|c|c|c|c|c|c|c|c|}
\hline
\textbf{Individuo} & \textbf{Educ} ($X_1$) & \textbf{Exp} ($X_2$)  & \textbf{Salario ($Y$)} & $\hat{Y}$ & $(Y_i - \Bar{Y})^2$ & $(\hat{Y_i} - \Bar{Y})^2$ & $(Y_i-\hat{Y_i})^2$\\
\hline
1 & 12 & 14 & 1000 & 990 & 57600 & 62500& \\
2 & 14 & 15 & 1200 & 1160 & 1600  & 6400 & \\
3 & 15 & 17 & 1300 & 1320 & 3600  & 6400& \\
4 & 16 & 15 & 1200 & 1280 & 1600  & 1600& \\
5 & 18 & 16 & 1500 & 1450 & 67600 & 44100& \\
\hline
\end{tabular}
\end{center}
SSE = 121000, por tanto SSR = 11000. 
$$R^2 = 1- \frac{11000}{132000}=0.92$$
y el ajustado
$$R^2_{adj}= 1- \frac{5-1}{5-2-1}\frac{11000}{132000}=0.83$$
    \item (5 puntos) Compare ambos modelos. 
    \\
    \textbf{Solución}
    \\
    Se puede apreciar que el $\beta_{educ}$ baja de valor pasando de 75 a 60, por lo que pesa menos al agregar experienia, por otra parte, el $R^2$ mejora considerablemente pasando de 0.85 a 0.92, y el $R^2$ ajustado de 0.8 a 0.83, por lo que se prodría decir que incorporando la experiencia se puede predecir de mejor manera el salario.
\end{enumerate}

\end{enumerate}



\end{enumerate}

\section{Formulario}


\begin{itemize}

\item  Estimadores MCO para RLS.
{
$$\hat{\beta}_0 = \bar{y} - \hat{\beta}_1 \bar{x} $$
$$\hat{\beta}_1 = \frac{\sum_{i=1}^{n}(x_i - \bar{x})(y_i-\bar{y})}{\sum_{i=1}^{n}(x_i - \bar{x})^2}=\frac{Cov(x,y)}{Var(x)}$$}


\item  Estimadores MCO para RLM.
{
$${\hat{\beta}} = (X^TX)^{-1} X^TY $$}


    \item \[ R^2  = 1 - \frac{SSR}{SST} \] 
\item  $$SST = \sum_{i=1}^{n} (y_i - \bar{y})^2 \quad \quad SSR = \sum_{i=1}^{n}  \hat{u}_i^2 $$
    \item \[ \bar{R}^2 =  1 - \frac{(n-1)SSR}{(n-k-1)SST} \]

    \item  $$\hat{\sigma}^2 = \frac{\sum_{i=1}^{n} \hat{u}_i^2 }{n-k-1}= \frac{SSR}{n - k - 1}$$
    \item \[
\text{Var}(\hat{\beta}|X) = \hat{\sigma}^2(X^{\top}X)^{-1}
\]
    \end{itemize}



\end{document}